\documentclass[10pt,a4paper]{amsart}
\usepackage{amsmath,amssymb,amsthm,graphicx}
\usepackage{verbatim}
\newenvironment{reference}{\comment}{\endcomment}
%\newenvironment{reference}{}{}
\newenvironment{slogan}{\comment}{\endcomment}
\newenvironment{history}{\comment}{\endcomment}

% For commutative diagrams we use Xy-pic
\usepackage[all]{xy}

% We use 2cell for 2-commutative diagrams.
\xyoption{2cell}
% Theorem environments.
%
\theoremstyle{plain}
\newtheorem{theorem}[subsection]{Theorem}
\newtheorem{proposition}[subsection]{Proposition}
\newtheorem{lemma}[subsection]{Lemma}

\theoremstyle{definition}
\newtheorem{definition}[subsection]{Definition}
\newtheorem{example}[subsection]{Example}
\newtheorem{exercise}[subsection]{Exercise}
\newtheorem{situation}[subsection]{Situation}

\theoremstyle{remark}
\newtheorem{remark}[subsection]{Remark}
\newtheorem{remarks}[subsection]{Remarks}

\numberwithin{equation}{subsection}

% Macros
%
\def\lim{\mathop{\mathrm{lim}}\nolimits}
\def\colim{\mathop{\mathrm{colim}}\nolimits}
\def\Spec{\mathop{\mathrm{Spec}}}
\def\Hom{\mathop{\mathrm{Hom}}\nolimits}
\def\Ext{\mathop{\mathrm{Ext}}\nolimits}
\def\SheafHom{\mathop{\mathcal{H}\!\mathit{om}}\nolimits}
\def\SheafExt{\mathop{\mathcal{E}\!\mathit{xt}}\nolimits}
\def\Sch{\mathit{Sch}}
\def\Mor{\mathop{\mathrm{Mor}}\nolimits}
\def\Ob{\mathop{\mathrm{Ob}}\nolimits}
\def\Sh{\mathop{\mathit{Sh}}\nolimits}
\def\NL{\mathop{N\!L}\nolimits}
\def\CH{\mathop{\mathrm{CH}}\nolimits}
\def\proetale{{pro\text{-}\acute{e}tale}}
\def\etale{{\acute{e}tale}}
\def\QCoh{\mathit{QCoh}}
\def\Ker{\mathop{\mathrm{Ker}}}
\def\Im{\mathop{\mathrm{Im}}}
\def\Coker{\mathop{\mathrm{Coker}}}
\def\Coim{\mathop{\mathrm{Coim}}}

% Boxtimes
%
\DeclareMathSymbol{\boxtimes}{\mathbin}{AMSa}{"02}

%
% Macros for moduli stacks/spaces
%
\def\QCohstack{\mathcal{QC}\!\mathit{oh}}
\def\Cohstack{\mathcal{C}\!\mathit{oh}}
\def\Spacesstack{\mathcal{S}\!\mathit{paces}}
\def\Quotfunctor{\mathrm{Quot}}
\def\Hilbfunctor{\mathrm{Hilb}}
\def\Curvesstack{\mathcal{C}\!\mathit{urves}}
\def\Polarizedstack{\mathcal{P}\!\mathit{olarized}}
\def\Complexesstack{\mathcal{C}\!\mathit{omplexes}}
% \Pic is the operator that assigns to X its picard group, usage \Pic(X)
% \Picardstack_{X/B} denotes the Picard stack of X over B
% \Picardfunctor_{X/B} denotes the Picard functor of X over B
\def\Pic{\mathop{\mathrm{Pic}}\nolimits}
\def\Picardstack{\mathcal{P}\!\mathit{ic}}
\def\Picardfunctor{\mathrm{Pic}}
\def\Deformationcategory{\mathcal{D}\!\mathit{ef}}
\usepackage{cleveref}
\crefname{originallemma}{original lemma}{original lemmas}
\Crefname{originallemma}{Original lemma}{Original lemmas}
\makeatletter
\newlabel{divisors-lemma-codim-1-part}{{0AGB}{source}{}{sourcepointer-1}{}}
\newlabel{divisors-lemma-codim-1-part@cref}{{[originallemma][0][]0AGB}{[1][0][]source}{}{}{}}
\makeatother
\crefname{originallemma}{original lemma}{original lemmas}
\Crefname{originallemma}{Original lemma}{Original lemmas}
\makeatletter
\newlabel{divisors-lemma-codimension-1-is-effective-Cartier}{{0BXH}{source}{}{sourcepointer-2}{}}
\newlabel{divisors-lemma-codimension-1-is-effective-Cartier@cref}{{[originallemma][0][]0BXH}{[1][0][]source}{}{}{}}
\makeatother
\crefname{originalequation}{original equation}{original equations}
\Crefname{originalequation}{Original equation}{Original equations}
\makeatletter
\newlabel{lemma-equation-plane-curve}{{0BYB}{source}{}{sourcepointer-3}{}}
\newlabel{lemma-equation-plane-curve@cref}{{[originalequation][0][]0BYB}{[1][0][]source}{}{}{}}
\makeatother
\title{A principal plane subscheme is a curve exactly when its equation is irreducible}
\author{The Stacks Project Authors}
\date{stacks-project:curves.tex:local-snapshot-20261012; GNU FDL 1.2 or later}
\begin{document}
\maketitle
\noindent\textit{Editorial note.} All delivered mathematical text is the author's own, copied from the pinned Stacks Project chapter file curves.tex; no statement or proof step was rewritten, reordered or removed.\par\medskip
\noindent\textit{Editorial note.} The counted proof is the author's complete argument: the reducible construction, the weighted decomposition with component equations, the scalar comparison of generators, and the prime-divisor conclusion are all retained verbatim.\par\medskip
\noindent\textit{Editorial note.} The label lemma-plane-curve names the counted statement, and the proof environment immediately following it in the source is delivered verbatim as its proof.\par\medskip
\noindent\textit{Editorial note.} The counted statement is carried by the author's own lemma environment in the chapter source.\par\medskip
\noindent\textit{Allowed standard prerequisites (editorial).}\begin{itemize}
\item The plane-curve setup. Same chapter, Lemma lemma-equation-plane-curve. Fixes the closed subscheme Z of the projective plane with principal ideal generated by the homogeneous form F.
\item Codimension-one divisors on the plane. Divisors, Lemmas divisors-lemma-codim-1-part and divisors-lemma-codimension-1-is-effective-Cartier. Provides the decomposition Z = sum a\_i D\_i into irreducible components with multiplicities, each D\_i cut out by a homogeneous F\_i.
\end{itemize}\medskip
\ifdefined\hypertarget\hypertarget{sourcepointer-1}{}\fi
\noindent\textit{Source reference.} Original lemma 0AGB: sources/stacks-chow/tags.txt, source lines 13189--13189. Prerequisite pointer invoked at proof line(s) 1892; it names a standard fact of a sibling chapter that the counted argument uses, and every step of the counted proof built on it is retained verbatim in the delivered body.\par\medskip
\ifdefined\hypertarget\hypertarget{sourcepointer-2}{}\fi
\noindent\textit{Source reference.} Original lemma 0BXH: sources/stacks-chow/tags.txt, source lines 14980--14980. Prerequisite pointer invoked at proof line(s) 1893; it names a standard fact of a sibling chapter that the counted argument uses, and every step of the counted proof built on it is retained verbatim in the delivered body.\par\medskip
\ifdefined\hypertarget\hypertarget{sourcepointer-3}{}\fi
\noindent\textit{Source reference.} Original equation 0BYB: sources/stacks-curves/pinned-original-source.tex, source lines 1859--1859. Prerequisite pointer invoked at proof line(s) 1880; it names a standard fact of the same chapter that the counted argument uses, and every step of the counted proof built on it is retained verbatim in the delivered body.\par\medskip
\begin{lemma}
\label{lemma-plane-curve}
Let $Z \subset \mathbf{P}^2_k$ be as in Lemma \ref{lemma-equation-plane-curve}
and let $I(Z) = (F)$ for some $F \in k[T_0, T_1, T_2]$.
Then $Z$ is a curve if and only if $F$ is irreducible.
\end{lemma}

\begin{proof}
If $F$ is reducible, say $F = F' F''$ then let $Z'$ be the closed subscheme
of $\mathbf{P}^2_k$ defined by $F'$. It is clear that $Z' \subset Z$
and that $Z' \not = Z$. Since $Z'$ has dimension $1$ as well, we conclude
that either $Z$ is not reduced, or that $Z$ is not irreducible.
Conversely, write $Z = \sum a_i D_i$ where $D_i$ are the irreducible
components of $Z$, see
Divisors, Lemmas \ref{divisors-lemma-codim-1-part} and
\ref{divisors-lemma-codimension-1-is-effective-Cartier}.
Let $F_i \in k[T_0, T_1, T_2]$ be the homogeneous
polynomial generating the ideal of $D_i$. Then it is clear that
$F$ and $\prod F_i^{a_i}$ cut out the same closed subscheme of
$\mathbf{P}^2_k$. Hence $F = \lambda \prod F_i^{a_i}$ for some
$\lambda \in k^*$ because both generate the ideal of $Z$.
Thus we see that if $F$ is irreducible, then $Z$ is
a prime divisor, i.e., a curve.
\end{proof}

\end{document}
